% Options for packages loaded elsewhere
% Options for packages loaded elsewhere
\PassOptionsToPackage{unicode}{hyperref}
\PassOptionsToPackage{hyphens}{url}
\PassOptionsToPackage{dvipsnames,svgnames,x11names}{xcolor}
%
\documentclass[
  letterpaper,
  DIV=11,
  numbers=noendperiod]{scrartcl}
\usepackage{xcolor}
\usepackage{amsmath,amssymb}
\setcounter{secnumdepth}{-\maxdimen} % remove section numbering
\usepackage{iftex}
\ifPDFTeX
  \usepackage[T1]{fontenc}
  \usepackage[utf8]{inputenc}
  \usepackage{textcomp} % provide euro and other symbols
\else % if luatex or xetex
  \usepackage{unicode-math} % this also loads fontspec
  \defaultfontfeatures{Scale=MatchLowercase}
  \defaultfontfeatures[\rmfamily]{Ligatures=TeX,Scale=1}
\fi
\usepackage{lmodern}
\ifPDFTeX\else
  % xetex/luatex font selection
\fi
% Use upquote if available, for straight quotes in verbatim environments
\IfFileExists{upquote.sty}{\usepackage{upquote}}{}
\IfFileExists{microtype.sty}{% use microtype if available
  \usepackage[]{microtype}
  \UseMicrotypeSet[protrusion]{basicmath} % disable protrusion for tt fonts
}{}
\makeatletter
\@ifundefined{KOMAClassName}{% if non-KOMA class
  \IfFileExists{parskip.sty}{%
    \usepackage{parskip}
  }{% else
    \setlength{\parindent}{0pt}
    \setlength{\parskip}{6pt plus 2pt minus 1pt}}
}{% if KOMA class
  \KOMAoptions{parskip=half}}
\makeatother
% Make \paragraph and \subparagraph free-standing
\makeatletter
\ifx\paragraph\undefined\else
  \let\oldparagraph\paragraph
  \renewcommand{\paragraph}{
    \@ifstar
      \xxxParagraphStar
      \xxxParagraphNoStar
  }
  \newcommand{\xxxParagraphStar}[1]{\oldparagraph*{#1}\mbox{}}
  \newcommand{\xxxParagraphNoStar}[1]{\oldparagraph{#1}\mbox{}}
\fi
\ifx\subparagraph\undefined\else
  \let\oldsubparagraph\subparagraph
  \renewcommand{\subparagraph}{
    \@ifstar
      \xxxSubParagraphStar
      \xxxSubParagraphNoStar
  }
  \newcommand{\xxxSubParagraphStar}[1]{\oldsubparagraph*{#1}\mbox{}}
  \newcommand{\xxxSubParagraphNoStar}[1]{\oldsubparagraph{#1}\mbox{}}
\fi
\makeatother


\usepackage{longtable,booktabs,array}
\newcounter{none} % for unnumbered tables
\usepackage{calc} % for calculating minipage widths
% Correct order of tables after \paragraph or \subparagraph
\usepackage{etoolbox}
\makeatletter
\patchcmd\longtable{\par}{\if@noskipsec\mbox{}\fi\par}{}{}
\makeatother
% Allow footnotes in longtable head/foot
\IfFileExists{footnotehyper.sty}{\usepackage{footnotehyper}}{\usepackage{footnote}}
\makesavenoteenv{longtable}
\usepackage{graphicx}
\makeatletter
\newsavebox\pandoc@box
\newcommand*\pandocbounded[1]{% scales image to fit in text height/width
  \sbox\pandoc@box{#1}%
  \Gscale@div\@tempa{\textheight}{\dimexpr\ht\pandoc@box+\dp\pandoc@box\relax}%
  \Gscale@div\@tempb{\linewidth}{\wd\pandoc@box}%
  \ifdim\@tempb\p@<\@tempa\p@\let\@tempa\@tempb\fi% select the smaller of both
  \ifdim\@tempa\p@<\p@\scalebox{\@tempa}{\usebox\pandoc@box}%
  \else\usebox{\pandoc@box}%
  \fi%
}
% Set default figure placement to htbp
\def\fps@figure{htbp}
\makeatother


% definitions for citeproc citations
\NewDocumentCommand\citeproctext{}{}
\NewDocumentCommand\citeproc{mm}{%
  \begingroup\def\citeproctext{#2}\cite{#1}\endgroup}
\makeatletter
 % allow citations to break across lines
 \let\@cite@ofmt\@firstofone
 % avoid brackets around text for \cite:
 \def\@biblabel#1{}
 \def\@cite#1#2{{#1\if@tempswa , #2\fi}}
\makeatother
\newlength{\cslhangindent}
\setlength{\cslhangindent}{1.5em}
\newlength{\csllabelwidth}
\setlength{\csllabelwidth}{3em}
\newenvironment{CSLReferences}[2] % #1 hanging-indent, #2 entry-spacing
 {\begin{list}{}{%
  \setlength{\itemindent}{0pt}
  \setlength{\leftmargin}{0pt}
  \setlength{\parsep}{0pt}
  % turn on hanging indent if param 1 is 1
  \ifodd #1
   \setlength{\leftmargin}{\cslhangindent}
   \setlength{\itemindent}{-1\cslhangindent}
  \fi
  % set entry spacing
  \setlength{\itemsep}{#2\baselineskip}}}
 {\end{list}}
\usepackage{calc}
\newcommand{\CSLBlock}[1]{\hfill\break\parbox[t]{\linewidth}{\strut\ignorespaces#1\strut}}
\newcommand{\CSLLeftMargin}[1]{\parbox[t]{\csllabelwidth}{\strut#1\strut}}
\newcommand{\CSLRightInline}[1]{\parbox[t]{\linewidth - \csllabelwidth}{\strut#1\strut}}
\newcommand{\CSLIndent}[1]{\hspace{\cslhangindent}#1}



\setlength{\emergencystretch}{3em} % prevent overfull lines

\providecommand{\tightlist}{%
  \setlength{\itemsep}{0pt}\setlength{\parskip}{0pt}}



 


\usepackage{booktabs}
\usepackage{caption}
\usepackage{longtable}
\usepackage{colortbl}
\usepackage{array}
\usepackage{anyfontsize}
\usepackage{multirow}
\KOMAoption{captions}{tableheading}
\makeatletter
\@ifpackageloaded{caption}{}{\usepackage{caption}}
\AtBeginDocument{%
\ifdefined\contentsname
  \renewcommand*\contentsname{Table of contents}
\else
  \newcommand\contentsname{Table of contents}
\fi
\ifdefined\listfigurename
  \renewcommand*\listfigurename{List of Figures}
\else
  \newcommand\listfigurename{List of Figures}
\fi
\ifdefined\listtablename
  \renewcommand*\listtablename{List of Tables}
\else
  \newcommand\listtablename{List of Tables}
\fi
\ifdefined\figurename
  \renewcommand*\figurename{Figure}
\else
  \newcommand\figurename{Figure}
\fi
\ifdefined\tablename
  \renewcommand*\tablename{Table}
\else
  \newcommand\tablename{Table}
\fi
}
\@ifpackageloaded{float}{}{\usepackage{float}}
\floatstyle{ruled}
\@ifundefined{c@chapter}{\newfloat{codelisting}{h}{lop}}{\newfloat{codelisting}{h}{lop}[chapter]}
\floatname{codelisting}{Listing}
\newcommand*\listoflistings{\listof{codelisting}{List of Listings}}
\makeatother
\makeatletter
\makeatother
\makeatletter
\@ifpackageloaded{caption}{}{\usepackage{caption}}
\@ifpackageloaded{subcaption}{}{\usepackage{subcaption}}
\makeatother
\usepackage{bookmark}
\IfFileExists{xurl.sty}{\usepackage{xurl}}{} % add URL line breaks if available
\urlstyle{same}
% fallback for those not using the hyperref driver hyperxmp:
\makeatletter
\@ifundefined{xmpquote}{\newcommand{\xmpquote}[1]{#1}}{}
\makeatother
\hypersetup{
  pdftitle={Replicate Boulesteix et al.'s simulation study},
  pdfauthor={Ludo Mels},
  colorlinks=true,
  linkcolor={blue},
  filecolor={Maroon},
  citecolor={Blue},
  urlcolor={Blue},
  pdfcreator={LaTeX via pandoc}}


\title{Replicate Boulesteix et al.'s simulation study}
\author{Ludo Mels}
\date{2026-11-10}
\begin{document}
\maketitle


\section{Introduction}\label{introduction}

This report aims to replicate the simulation study by Boulesteix et al.
(2020). This study includes a small simulation study that investigated
the impact of measurement error on the results. In order to this this,
they used real medical data and assumed these data included no
measurement error. Subsequently, they added measurement error to both
the exposure variable (\texttt{HbA1C}) and the confounder variable
(\texttt{bmi}), and refitted the model. While the supplementary
materials of this study included an \texttt{R}-script, the report still
lacked reproducibility. For example, it did not include information on
the computational environment.

As part of the course \href{https://mlarpis.github.io/markup/}{Markup
Languages and Reproducible Programming in Statistics}, this report aims
to improve the reproducability of the study by Boulesteix et al. (2020).
The method's section discusses the changes made to the code syntax and
the computational environment to improve the reproducibility of the
report. The results provides a textual description and an interpretation
of the simulation study results.

\section{Method}\label{method}

The core aim of this report is to improve the reproducibility of the
\texttt{R} script as provided by Boulesteix et al. (2020). First, the
required packages were loaded and registered in the \texttt{renv.lock}
file using the \texttt{renv} package (Ushey \& Wickham, 2026). Note
that, since this project is part of the overall
\texttt{markup\_course}-repository, a base path is defined to direct the
paths towards the subfolders. Also note that the reproducible
environment (i.e., the \texttt{renv.lock} file) and the Quarto project
are defined at the root of the project.

Besides the lack of information about the computational environment, the
code provided in the supplementary materials (Boulesteix et al., 2020)
lacked readability. Several improvements have been made in terms of
readability. First, pipe operators (i.e.,
\texttt{\textbar{}\textgreater{}}) were introduced in appropriate
places. For example, only a subset of variables were used. Instead of
first loading the full data and then selecting these relevant variables,
the variables were immediately extracted using the pipe operator.
Second, some objects were renamed to be more informative. For example,
the object \texttt{d} was changed to \texttt{dat\_allcases}, and
\texttt{dc} (i.e., the data including the complete cases only) was
changed to \texttt{dat}, since this would be the final data used in the
analyses.

In addition to the improving the readability of the code, the efficiency
was also improved. Note that the improvement of the code's efficiency
often also results in improved readability. The following improvements
have been made. First, instead of using multiple lines for merging, the
data were stored as a list object and subsequently merged in one line to
become one data frame. Also, the variables were renamed using the
\texttt{names()}-function, instead of creating copies in the same data
frame with different names. Thirdly, the old script used to perform
\texttt{dat\$age{[}d\$age\textless{}18{]}\ \textless{}-\ NA}. Since
there were no cases with age below 18, this line has been removed. Also
note that that line of code only removes the age for these respondents
and not their complete entry, which is remarkable.

In some cases, extra objects were defined to improve efficieny and
readability. For example, the number of scenario's was used in multiple
places in the script, but it had never been defined. Instead of
hard-coding \texttt{dim(scenarios){[}1{]}} each time, this script
defined \texttt{n.scenarios} first. Also, the original models were
fitted twice; one time for extracting the model summary and one time for
extracting the confidence interval. In the new script, the model is
stored in an object and the model summary and confidence intervals were
requested afterwards. In another place, the number of observations
\texttt{n} was defined as an object, although it was not used in the
following syntax (i.e., again, the number of observations was hard-coded
in these following lines).

converted sex to factor outside the models

The for-loop used in the simulation study was partially redesigned.
First of all, the seed was included in the loop. Because the seed was
repeated at each iteration, the randomness of the process was disturbed.
Second, to meet conventions, \texttt{i} and \texttt{j} were used as
indices for the rows and columns of the results-matrix, respectively.
Finally, instead of printing the iteration number at the start of each
loop, the progress was mentioned as a pecentage for each 100 iterations.

Finally, to serve the results section, the model coefficients were
stored as an object. This allows for direct implementation in the text.
Also, instead of displaying the absolute values of the regression
coeffients in the figure, the figure now displays the difference to the
original coefficient. Although the colors in the original heat-map could
already be interpreted as the difference to the original, it can be
difficult to see exactly in which situations the difference was
greatest.

\section{Results}\label{results}

\begin{table}

\caption{\label{tbl-model_summary}Summary of the original model}

\centering{

\fontsize{12.0pt}{14.0pt}\selectfont
\begin{tabular*}{\linewidth}{@{\extracolsep{\fill}}lccc}
\toprule
\textbf{Characteristic} & \textbf{Beta} & \textbf{95\% CI} & \textbf{p-value} \\ 
\midrule\addlinespace[2.5pt]
Glycohemoglobin (\%) & 0.74 & 0.36, 1.1 & {\bfseries <0.001} \\ 
Body Mass Index (kg/m**2) & 0.31 & 0.25, 0.37 & {\bfseries <0.001} \\ 
Age in years at screening & 0.43 & 0.41, 0.45 & {\bfseries <0.001} \\ 
sex &  &  &  \\ 
    1 & — & — &  \\ 
    2 & -3.8 & -4.5, -3.0 & {\bfseries <0.001} \\ 
\bottomrule
\end{tabular*}
\begin{minipage}{\linewidth}
\vspace{.05em}
Abbreviation: CI = Confidence Interval\\
\end{minipage}

}

\end{table}%

The original linear regression model is summarized in
Table~\ref{tbl-model_summary}. In the simulation study, this model is
assumed to have no measurement error. In the original model, the
regression coeffient for the exposure variable \texttt{HbA1C} was 0.736,
and the coefficient for the confounder variable was 0.31. Both variables
were significant at the level of \(p < .001\).

In the simulation study, measurement error was artificially added to the
raw data. The results of the simulation study are summarized in
Figure~\ref{fig-simulation_res}. The percentage of injected measurement
error to the exposure and confounder variables are displayed on the x
and y-axis, respectively. The contents of the figure reflect the
averaged difference to the actual model coefficient for the exposure
variable \texttt{HbA1C}.

The plot shows that the estimates are most impacted by measurement error
in the exposure variable. Remarkably, if measurement error in the
exposure variable is combined with measurement error in the confounder
variable, the difference to the model's actual coeffient descreases.
This can be explained by the fact that, while the estimated coefficient
decreases when measurment error is added to the exposure variable, it
actually increases when measurement error is added to the confounder
variable. Thus, in situarions where variables interact (e.g., when a
confounder is present), more measurement error does not necessarily
result in a more biased estimate.

\begin{figure}

\centering{

\pandocbounded{\includegraphics[keepaspectratio]{report_files/figure-pdf/fig-simulation_res-1.pdf}}

}

\caption{\label{fig-simulation_res}}

\end{figure}%

\section{Conclusions}\label{conclusions}

The goal of this report was to improve the replicability of the
simulation study conducted by Boulesteix et al. (2020). The code
provided in the original study was improved in terms of readability and
efficiency. In addition, a reproducible environment was created using
the \texttt{renv} package (Ushey \& Wickham, 2026).

The results were similar to the results provided by Boulesteix et al.
(2020) in substantive terms. Nevertheless, the exact estimates differed
because the seed-setting was changed in the code. The original seed
setting disturbed the randomness by repeating the seed-setting in the
loop.

\protect\phantomsection\label{refs}
\begin{CSLReferences}{1}{0}
\bibitem[\citeproctext]{ref-Boulesteix2020}
Boulesteix, A.-L., Groenwold, R. H., Abrahamowicz, M., Binder, H.,
Briel, M., Hornung, R., Morris, T. P., Rahnenführer, J., \& Sauerbrei,
W. (2020). Introduction to statistical simulations in health research.
\emph{BMJ Open}, \emph{10}(12).
\url{https://doi.org/10.1136/bmjopen-2020-039921}

\bibitem[\citeproctext]{ref-renv}
Ushey, K., \& Wickham, H. (2026). \emph{Renv: Project environments}.
\url{https://doi.org/10.32614/CRAN.package.renv}

\end{CSLReferences}




\end{document}
